\documentclass{chsmc}
\chsmcsetup{
	year = 2010,
	part = 1,
	type = A,
	show-answers = false,
}
\begin{document}

\section{填空题：本大题共 8 小题，每小题 8 分，共 64 分.}

\begin{enumerate}
	\item 函数 $f(x) = \sqrt{x-5} - \sqrt{24-3x}$ 的值域是 \blankline
	\item 已知函数 $y = (a\cos^2 x - 3)\sin x$ 的最小值为 $-3$，
		则实数 $a$ 的取值范围是 \blankline
	\item 双曲线 $x^2-y^2=1$ 的右半支与直线 $x=100$ 围成的区域内部
		（不含边界）整点（纵横坐标均为整数的点）的个数是 \blankline
	\item 已知 $\{a_n\}$ 是公差不为 $0$ 的等差数列，$\{b_n\}$ 是等比数列，
		其中 $a_1 = 3$, $b_1 = 1$, $a_2 = b_2$, $3a_5 = b_3$，且存在常数 $\alpha$, $\beta$
		使得对每一个正整数 $n$ 都有 $a_n = \log_{\alpha} b_n + \beta$，
		则 $\alpha + \beta =$ \blankline
	\item 函数 $f(x) = a^{2x} + 3a^x - 2$（$a > 0$, $a \neq 1$）在区间
		$x\in[-1,1]$ 上的最大值为 $8$，则它在这个区间上的最小值是 \blankline
	\item 两人轮流投掷骰子，每人每次投掷两颗，第一个使两颗骰子点数和大于 $6$ 者为胜，
		否则轮由另一人投掷. 先投掷人的获胜概率是 \blankline
	\item 正三棱柱 $ABC$-$A_1B_1C_1$ 的 $9$ 条棱长都相等，$P$ 是 $CC_1$ 的中点，
		二面角 $B$-$A_1P$-$B_1 = \alpha$，则 $\sin\alpha =$ \blankline
	\item 方程 $x + y + z = 2010$ 满足 $x \leqslant y \leqslant z$ 的正整数解
		$(x,y,z)$ 的个数是 \blankline
\end{enumerate}

\section{解答题：本大题共 3 个小题，满分 56 分. 解答应写出文字说明、证明过程或演算步骤.}

\begin{enumerate}[resume]
	\item (本题满分 16 分) 已知函数 $f(x) = ax^3 + bx^2 + cx + d$（$a \neq 0$），
		当 $0 \leqslant x \leqslant 1$ 时，$|f'(x)| \leqslant 1$，试求 $a$ 的最大值.
	\item (本题满分 20 分) 已知抛物线 $y^2 = 6x$ 上的两个动点
		$A(x_1,y_1)$ 和 $B(x_2,y_2)$，其中 $x_1 \neq x_2$，且 $x_1 + x_2 = 4$.
		线段 $AB$ 的垂直平分线与 $x$ 轴交于点 $C$，求 $\triangle ABC$ 面积的最大值.
	\item (本题满分 20 分) 证明：方程 $2x^3 + 5x - 2 = 0$ 恰有一个实数根 $r$，
		且存在唯一的严格递增正整数数列 $\{a_n\}$，使得
		$\dfrac{2}{5} = r^{a_1} + r^{a_2} + r^{a_3} + \cdots$.
\end{enumerate}

\end{document}
